\section{State of the Art on Question Answering with External Knowledge}
\label{sec:sota:qa}

% from refact dog gcr
\todo{improve}
\begin{definition}[Question Answering]
    \emph{\Acf{qa}} refers to the task of producing a concise and relevant answer to a question expressed in natural language~\cite{entityorientedsearch2018balog}.
\end{definition}
\Ac{qa} systems may be focused on a specific domain, \eg supporting only questions regarding law---in this cas we speak of \emph{closed-domain \acl{qa}}, or designed to support broad domain-free questions, \eg questions referring to large general-purpose \ac{kr}, such as Wikipedia~\cite{zaib2022conversationalqasurvey}. As a consequence, open-domain \ac{qa} systems tend to more generalizable~\cite{zaib2022conversationalqasurvey}, thus, in this thesis we concentrate on open-domain \ac{qa}.

These system must acquire external information to answer the given question and this information is generally acquired from unstructured textual document collections, structured \aclp{kb}, or a combination of both~\cite{caballero2021brief}. Modern \acp{llm} contain considerable factual knowledge in the parameters~\cite{chang_how_do_llm_acquire_factual_knowledge_2024,physics_of_lm_3.1_2024,petroni-etal-2019-language-models-kbs,roberts-etal-2020-how-much-kg-in-llm}, but information is constantly evolving and updating \ac{llm} knowledge generally requires expensive fine-tuning~\cite{kg_editing_survey_2024}.

% retrieval reader

Two main paradigms exist for injecting external knowledge. The first is \emph{input-based KE-QA},
where the model is provided with external knowledge as input (plain text or dense vectors).
This is instantiated by \acf{rag} methods~\cite{gao2024retrievalaugmentedgenerationlargelanguage,karpukhin-etal-2020-dense,lewis2020retrievalaugmented},
including agent-style prompting that calls tools at inference time~\cite{yao2023reactsynergizingreasoningacting,NEURIPS2023_d842425e}.
Such methods reduce hallucination~\cite{shuster-etal-2021-retrieval-augmentation} but increase latency and complexity,
and inflate the number of input tokens~\cite{yang2024memory3,sun2025zerosearchincentivizesearchcapability}.
The second is \emph{memory-based KE-QA}, where architectures are equipped with non-parametric memory
modules coupled to the LLM via cross-attention or learned keys~\cite{wu-etal-2022-efficient,yang2024memory3}.
These integrate knowledge more tightly but require architectural modifications, reducing plug-and-play
applicability of existing pretrained LLMs.

\todo{go on}

\paragraph{KGQA as a subcase of knowledge-enhanced QA.}
Knowledge Graph QA (KGQA) historically pursued access to external knowledge through \emph{executable
queries} on a graph store. Semantic-parsing approaches~\cite{Lan_Jiang_2020,Ye_2022,Yu_2023} translate questions
into logical queries (e.g.\ SPARQL) and execute them against the KG. This can be viewed as a
KE-QA instance where KG access is \emph{delegated to a query compiler} rather than retrieved textually.
However, such pipelines require gold logical forms for training and may generate non-executable queries.

More recently, LLM-based KGQA~\cite{Salnikov_2023,Wu_2023,He_2024} has aligned with the two KE-QA paradigms above:
(i) input-based KGQA, where KG facts are first retrieved (explicitly or via learned retrievers) and fed to
an LLM~\cite{Jin_2024,Sun_2024b,Ren_2024,Luo_2024,Mavromatis_Karypis_2024,Zhuang_2024}; and (ii) interaction-based KGQA, where LLMs
are prompted to traverse the KG step by step (e.g.\ ToG~\cite{Sun_2024a} or Tree-of-Traversals~\cite{Markowitz_2024}),
thus using the KG as an external structured search space. Both classes ultimately turn KG knowledge
into textual input to an LLM, thus inheriting input-growth and latency issues.

\paragraph{Constrained generation for KE-QA.}
A third, more recent direction incorporates external knowledge not by feeding or storing it, but by
\emph{constraining the search space of generation}~\cite{cao2021autoregressive,scholak-etal-2021-picard,bevilacqua2022autoregressive,geng-etal-2023-grammar}.
Constrained decoding enforces that output tokens conform to a target grammar or vocabulary space,
without injecting external text. When applied to KE-QA with KGs~\cite{li2024decodinggraphsfaithfulsound,luo2024graphconstrainedreasoningfaithfulreasoning},
the model is guided to generate KG-consistent triples or paths, grounding reasoning in the graph while
avoiding explicit retrieval. Existing constrained KE-QA approaches, however, either assume access to
small KGs or still depend on entity linking for grounding, thus retaining some limitations of input-based
pipelines. Our work positions itself in this constrained generation line, using constrained decoding
directly over large-scale KG facts (Wikidata-scale) without relying on retrieval or architectural changes.
